% ============================================================
%  Cu₃O₂ Kagome Monolayer on LaAlO₃(111)
%  Growth and Electrical Characterization Specification
%  For contract/facility use — no theoretical framework
% ============================================================
\documentclass[11pt, a4paper]{article}

\usepackage[margin=2.2cm]{geometry}
\usepackage{amsmath, amssymb}
\usepackage{booktabs, array}
\usepackage{xcolor}
\usepackage{hyperref}
\usepackage{enumitem}
\usepackage{fancyhdr}
\usepackage{setspace}

\hypersetup{colorlinks=true, linkcolor=blue!50!black, urlcolor=blue!60!black}
\definecolor{nm-gray}{RGB}{100,100,100}
\definecolor{nm-green}{RGB}{20,100,40}
\definecolor{nm-red}{RGB}{160,30,30}

\pagestyle{fancy}\fancyhf{}
\rhead{\textcolor{nm-gray}{\small Cu$_3$O$_2$ Growth \& Test Spec}}
\lhead{\textcolor{nm-gray}{\small Joseph Forrester}}
\cfoot{\thepage}

\setlength{\parindent}{0pt}\setlength{\parskip}{6pt}\onehalfspacing

\begin{document}

\begin{center}
  {\LARGE \textbf{Cu$_3$O$_2$ Kagome Monolayer}}\\[4pt]
  {\LARGE \textbf{on LaAlO$_3$(111)}}\\[6pt]
  {\large Growth and Electrical Characterization Specification}\\[10pt]
  {\color{nm-gray}\small Joseph Forrester \quad|\quad
    Independent Researcher}\\[4pt]
  {\color{nm-gray}\small May 2026 \quad|\quad Revised from April 2026}
\end{center}

\vspace{1em}\hrule\vspace{1em}

% ============================================================
\section*{Overview}

This document specifies the growth of a Cu$_3$O$_2$ sub-monolayer
on LaAlO$_3$(111) by molecular beam epitaxy, followed by
electrostatic gating and electrical transport measurement.
The deliverables are:

\begin{enumerate}[nosep]
  \item Verified Cu$_3$O$_2$ kagome monolayer on LAO(111)
        (RHEED, XPS confirmation)
  \item Sheet resistance vs.\ temperature ($R_\square(T)$)
        from 400~K to 2~K at multiple gate voltages
  \item Raw data files for all measurements
\end{enumerate}

% ============================================================
\section{Substrate}

\begin{center}
\begin{tabular}{@{}lp{9cm}@{}}
\toprule
\textbf{Parameter} & \textbf{Specification} \\
\midrule
Material & LaAlO$_3$ (LAO), single crystal --- primary \\
 & SrTiO$_3$ (STO), single crystal --- secondary
  (1--2 samples for strain comparison) \\
Orientation & (111) face (both substrates) \\
Miscut & $< 0.3^\circ$ \\
Size & 5~mm $\times$ 5~mm or 10~mm $\times$ 10~mm \\
Quantity & 5 LAO + 2 STO substrates minimum \\
Supplier & MTI Corporation, CrysTec, or equivalent \\
\bottomrule
\end{tabular}
\end{center}

% ============================================================
\section{Substrate Preparation}

\begin{center}
\begin{tabular}{@{}lp{9cm}@{}}
\toprule
\textbf{Step} & \textbf{Specification} \\
\midrule
1. Solvent clean & Ultrasonic in acetone (5~min), isopropanol
  (5~min), deionized water rinse, N$_2$ dry \\
2. Anneal & $1000^\circ$C in $10^{-6}$~Torr O$_2$ for 1~hour \\
3. Verification & RHEED: sharp $(1 \times 1)$ streaks confirming
  atomically flat, ordered surface \\
\bottomrule
\end{tabular}
\end{center}

% ============================================================
\section{Cu$_3$O$_2$ Deposition}

\begin{center}
\begin{tabular}{@{}lp{9cm}@{}}
\toprule
\textbf{Parameter} & \textbf{Specification} \\
\midrule
Source & E-beam evaporator, 99.999\% Cu \\
Substrate temperature & $650^\circ$C ($\pm 10^\circ$C) \\
Oxygen source & RF oxygen plasma (remote configuration),
  equivalent flux to $5 \times 10^{-5}$~Torr molecular O$_2$ \\
Deposition rate & 0.005~\AA/s (monitored by quartz crystal
  microbalance) \\
Target coverage & 0.5 monolayer equivalent \\
Estimated duration & $\sim 25$~minutes \\
\bottomrule
\end{tabular}
\end{center}

\textbf{In-situ monitoring:}
\begin{itemize}[nosep]
  \item RHEED: observe transition from $(1 \times 1)$ to
        $(2 \times 2)$ or
        $(\sqrt{3} \times \sqrt{3})R30^\circ$ reconstruction.
        Either pattern confirms kagome ordering on the (111)
        surface.
  \item If RHEED shows 3D spots or no reconstruction, the growth
        has failed. Adjust temperature or O$_2$ pressure.
\end{itemize}

\textbf{Post-deposition:}
\begin{itemize}[nosep]
  \item Optional post-anneal: $650^\circ$C,
        $5 \times 10^{-5}$~Torr O$_2$, 10~minutes
  \item Cool to room temperature under O$_2$ to maintain
        Cu$^{2+}$ oxidation state
\end{itemize}

\textbf{Why these conditions:}
\begin{itemize}[nosep]
  \item $650^\circ$C provides sufficient Cu surface diffusion
        for thermodynamic equilibrium while remaining below the
        Cu desorption threshold. Window: 400--750$^\circ$C.
  \item RF oxygen plasma (remote): Atomic oxygen oxidizes Cu
        to Cu$^{2+}$ on contact, reducing the metallic Cu$^0$
        lifetime by $\sim 1000\times$ compared to molecular
        O$_2$. This suppresses Cu-Cu clustering and favors
        2D kagome formation. Remote plasma configuration
        prevents energetic ion damage to the substrate.
        Equivalent O flux to $5 \times 10^{-5}$~Torr
        molecular O$_2$.
  \item 0.005~\AA/s ensures each deposited atom has time
        ($\sim 10^{11}$ hops) to reach its equilibrium position
        before the next atom arrives. This prevents kinetic
        trapping.
  \item 0.5~ML prevents 3D island formation while providing
        sufficient coverage for a connected kagome network.
\end{itemize}

% ============================================================
\section{Post-Growth Verification}

\begin{center}
\begin{tabular}{@{}llp{6cm}@{}}
\toprule
\textbf{Method} & \textbf{Target} & \textbf{Accept / Reject} \\
\midrule
RHEED & $(2 \times 2)$ streaks &
  Accept: streaky reconstruction.
  Reject: 3D spots or $(1 \times 1)$ only. \\
XPS & Cu 2p$_{3/2}$ &
  Accept: $933.5 \pm 0.5$~eV (Cu$^{2+}$).
  Reject: $932.5$~eV (Cu$^0$ metal) or
  $935$~eV (Cu$^{3+}$). \\
STM (if available) & Surface structure &
  Accept: corner-sharing triangles with
  $5.4 \pm 0.5$~\AA\ periodicity.
  Reject: disordered or island morphology. \\
\bottomrule
\end{tabular}
\end{center}

\textbf{Note:} STM is desirable but not required for the
transport measurement to proceed. RHEED + XPS are sufficient
to confirm the correct phase.

% ============================================================
\section{Electrical Contact Fabrication}

Four-point probe geometry for sheet resistance measurement:

\begin{itemize}[nosep]
  \item Deposit 4 contact pads (Ti/Au, 5/50~nm) by e-beam
        evaporation through a shadow mask, in van der Pauw
        configuration (4 corners of the sample)
  \item Wire bond or use indium press contacts to measurement
        leads
  \item Verify ohmic contact at room temperature before
        proceeding
\end{itemize}

% ============================================================
\section{Electrostatic Gating}

\begin{center}
\begin{tabular}{@{}lp{9cm}@{}}
\toprule
\textbf{Parameter} & \textbf{Specification} \\
\midrule
Ionic liquid & DEME-TFSI
  (N,N-diethyl-N-methyl-N-(2-methoxyethyl)ammonium
  bis(trifluoromethylsulfonyl)imide) \\
Application & Drop of ionic liquid covering the film surface
  between the contact pads \\
Gate electrode & Pt wire or coil suspended in the ionic liquid,
  not touching the film \\
Gate voltage range & $V_g = 0$ to $+4.5$~V in 0.5~V steps \\
Leakage check & Gate current $< 1$~$\mu$A at each $V_g$ \\
\bottomrule
\end{tabular}
\end{center}

\textbf{Note:} Ionic liquid gating must be applied at or above
room temperature (the ionic liquid freezes below $\sim 220$~K).
Set the gate voltage at room temperature, allow 5--10 minutes
for the electric double layer to form, then cool while
maintaining $V_g$.

% ============================================================
\section{Transport Measurement}

\subsection{Primary Measurement: $R_\square(T)$}

\begin{center}
\begin{tabular}{@{}lp{9cm}@{}}
\toprule
\textbf{Parameter} & \textbf{Specification} \\
\midrule
Instrument & PPMS, cryostat, or any system capable of
  controlled temperature sweep with electrical measurement \\
Temperature range & 400~K down to 2~K \\
Sweep direction & \textbf{Cool from 400~K to 2~K}
  (start above 400~K) \\
Sweep rate & 1--2~K/min \\
Measurement current & 1--10~$\mu$A (AC or DC, low enough to
  avoid heating the monolayer) \\
Data points & Record $R_\square$ at every 0.5~K or continuous
  logging \\
\bottomrule
\end{tabular}
\end{center}

\subsection{Measurement Sequence}

Perform the following runs on each verified sample:

\begin{enumerate}[nosep]
  \item $V_g = 0$~V: baseline (film likely insulating)
  \item $V_g = +1.0$~V: light gating
  \item $V_g = +2.0$~V
  \item $V_g = +3.0$~V
  \item $V_g = +4.0$~V
  \item $V_g = +4.5$~V: maximum gating
\end{enumerate}

At each gate voltage:
\begin{itemize}[nosep]
  \item Set $V_g$ at 300~K, wait 10~minutes
  \item Heat to 400~K
  \item Cool to 2~K at 1--2~K/min while logging $R_\square(T)$
  \item Warm back to 300~K before changing $V_g$
\end{itemize}

\subsection{What to Look For}

\begin{itemize}[nosep]
  \item A sharp drop in $R_\square$ at any temperature during
        cooling indicates a phase transition
  \item If $R_\square$ drops to the instrument noise floor
        ($< 0.01~\Omega$/sq), this indicates zero-resistance
        state
  \item Record the onset temperature of any transition as
        $T_{\text{onset}}$ and the zero-resistance temperature
        as $T_{\text{zero}}$
  \item If no transition is observed at any gate voltage,
        record the $R_\square(T)$ curves as-is --- this is
        still useful data
\end{itemize}

% ============================================================
\section{Deliverables}

\begin{enumerate}[nosep]
  \item RHEED images: pre-deposition and post-deposition
  \item XPS spectrum of Cu 2p region
  \item STM image (if available)
  \item $R_\square(T)$ data for each gate voltage (raw data
        files, not just plots)
  \item Gate leakage current vs.\ $V_g$ at room temperature
  \item Photos of sample with contacts and ionic liquid
        applied
  \item Any anomalies or deviations from the specified
        protocol, documented
\end{enumerate}

% ============================================================
\section{Sample Handling Notes}

\begin{itemize}[nosep]
  \item Store samples in vacuum or dry N$_2$ after growth
        (Cu$^{2+}$ may reduce in ambient over days/weeks)
  \item Handle with non-metallic tweezers only (avoid
        contamination)
  \item The ionic liquid is hygroscopic --- apply in a
        glovebox or dry environment
  \item Samples can be shipped in vacuum-sealed containers
        if growth and measurement are at different facilities
\end{itemize}

% ============================================================
\section{Interface Chemistry (for grower reference)}

Cu does not react with or diffuse into LaAlO$_3$:

\begin{center}
\begin{tabular}{@{}lrc@{}}
\toprule
\textbf{Reaction} & $\Delta H$ (kJ/mol) & \textbf{Occurs?} \\
\midrule
Cu$_3$O$_2$ on LAO surface & $-365$ &
  \textcolor{nm-green}{\textbf{Yes (favorable)}} \\
Cu substitutes for Al in LAO & $+694$ &
  \textcolor{nm-red}{No (forbidden)} \\
La$_2$CuO$_4$ at interface & $+144$ &
  \textcolor{nm-red}{No (forbidden)} \\
\bottomrule
\end{tabular}
\end{center}

Al diffusion from LAO at $650^\circ$C: $< 0.01$~\AA\ in
2~hours. The interface remains atomically sharp.

\vspace{2em}\hrule\vspace{0.5em}
{\color{nm-gray}\small Joseph Forrester --- Independent
Researcher --- May 2026 (Rev.\ 2)}

\end{document}
